\subsection{特征 0 时余代数 U(g) 的结构}

本号中，将假设 K 是特征 0 的域。

设 \(S(\mathfrak{g})\) 为向量空间 \(\mathfrak{g}\) 的对称代数，\(c_{\mathrm{s}}\) 为其余积（《代数》，第 III 章，第 11 节，第 1 号，例 6），\(\eta\) 为向量空间 \(\mathsf{S}(\mathfrak{g})\) 到向量空间 U 上的典范同构（第 I 章，第 2 节，第 7 号）。回忆若 \(x_{1},\ldots ,x_{n}\) 属于 \(\mathfrak{g}\)，则

\[
\eta (x _ {1} \dots x _ {n}) = \frac {1}{n !} \sum_ {\tau \in \mathfrak {S} _ {n}} \sigma (x _ {\tau (1)}) \dots \sigma (x _ {\tau (n)}).\tag{9}
\]

特别地，对 \(x \in g\) 与 \(n \geqslant 0\)，

\[
\eta (x ^ {n}) = \sigma (x) ^ {n}.\tag{10}
\]

注意，由《代数》第 III 章第 6 节第 1 号注 3，\(\eta\) 是 \(\mathbb{S}(\mathfrak{g})\) 到 U 的满足条件 (10) 的唯一线性映射。

命题 9. 对每个整数 \(n \geqslant 0\)，设 \(\mathbf{U}^n\) 为 \(\mathbf{U}\) 的由诸 \(\sigma(x)^n\)（\(x \in \mathfrak{g}\)）生成的向量子空间。

\begin{enumerate}
\def\labelenumi{(\alph{enumi})}
\tightlist
\item
  序列 \((\mathbf{U}^n)_{n\geqslant 0}\) 是向量空间 \(\mathbf{U}\) 的与其余代数结构相容的分次。
\end{enumerate}

给 U 赋以分次 (U \(^{n}\))。

\begin{enumerate}
\def\labelenumi{(\alph{enumi})}
\setcounter{enumi}{1}
\tightlist
\item
  典范映射 \(\eta: \mathbf{S}(\mathfrak{g}) \to \mathbf{U}\) 是分次余代数的同构。设 \(x \in \mathfrak{g}\)，\(n \in \mathbf{N}\)。那么
\end{enumerate}

\[
c _ {\mathrm{S}} (x ^ {n}) = c _ {\mathrm{S}} (x) ^ {n} = (x \otimes 1 + 1 \otimes x) ^ {n} = \sum_ {i = 0} ^ {n} {\binom {n} {i}} x ^ {i} \otimes x ^ {n - i}\tag{11}
\]

因为 \(c_{\mathbb{S}}\) 是代数同态。类似地，由 (10) 得

\[
\begin{array}{l} c (\eta (x ^ {n})) = c (\sigma (x) ^ {n}) = c (\sigma (x)) ^ {n} = (\sigma (x) \otimes 1 + 1 \otimes \sigma (x)) ^ {n} \\ = \sum_ {i = 0} ^ {n} \binom {n} {i} \sigma (x) ^ {i} \otimes \sigma (x) ^ {n - i} = \sum_ {i = 0} ^ {n} \binom {n} {i} \eta (x ^ {i}) \otimes \eta (x ^ {n - i}), \end{array}\tag{12}
\]

由此

\[
(\eta \otimes \eta) (c _ {\mathrm{S}} (x ^ {n})) = c (\eta (x ^ {n})).
\]

由于诸 \(x^n\)（对 \(x \in \mathfrak{g}\) 与 \(n \in \mathbf{N}\)）生成向量空间 \(\mathsf{S}(\mathfrak{g})\)，故 \((\eta \otimes \eta) \circ c_{\mathbf{S}} = c \circ \eta\)，从而 \(\eta\) 是余代数同构。

另一方面，公式 (10) 表明 \(\eta(\mathbb{S}^n(\mathfrak{g})) = \mathrm{U}^n\)，注意到 \(\mathbb{S}(\mathfrak{g})\) 的分次与其余代数结构相容这一事实，这就完成了 (a) 与 (b) 的证明。

U 的分次 \((\mathrm{U}^{n})_{n\geqslant0}\) 称为典范分次。

推论. 典范映射 \(\sigma\) 给出 \(g\) 到李代数 \(\mathrm{P}(\mathbf{U})\)（\(\mathbf{U}\) 的本原元素组成的李代数）上的同构。

由于 \(c^{+}\) 是次数 0 的分次同态，

\[
\mathrm{P} (\mathrm{U}) = \sum_ {n \geqslant 1} (\mathrm{P} (\mathrm{U}) \cap \mathrm{U} ^ {n}).
\]

只需证明：若 \(n > 1\) 且 \(a \in \mathrm{U}^n\) 是本原的，则 \(a = 0\)。今 \(a\) 可写成 \(\sum_{i} \lambda_i a_i^n\)，其中 \(\lambda_i \in \mathbf{K}\)，\(a_i \in \sigma(g)\)。由 (12)，双次数 \((1, n - 1)\) 的项在 \(c^+(a)\) 中为 \(n \sum_{i} \lambda_i a_i \otimes a_i^{n-1}\)。于是 \(\sum_{i} \lambda_i a_i \otimes a_i^{n-1} = 0\)。若 \(\mu: \mathrm{U} \otimes \mathrm{U} \to \mathrm{U}\) 是由 \(\mathrm{U}\) 上乘法所定义的线性映射，则

\[
a = \sum_ {i} \lambda_ {i} a _ {i} ^ {n} = \mu \left(\sum_ {i} \lambda_ {i} a _ {i} \otimes a _ {i} ^ {n - 1}\right) = 0.
\]

注. (1) \(\mathrm{U}_n = \sum_{i=0}^{n} \mathrm{U}^i\)（第 I 章，第 2 节，第 7 号，定理 1 的推论 4）。

\begin{enumerate}
\def\labelenumi{(\arabic{enumi})}
\setcounter{enumi}{1}
\tightlist
\item
  映射 \(\eta\) 是 \(\mathsf{S}(\mathfrak{g})\) 到 U 的分次余代数态射中使得 \(\eta(1) = 1\) 且 \(\eta(x) = \sigma(x)\)（对 \(x \in \mathfrak{g}\)）的唯一态射。因为若 \(\eta'\) 是满足这些条件的态射，我们对 \(n\) 用归纳法证明：\(\eta'(x^n) = \eta(x^n)\)（对 \(x \in \mathfrak{g}\) 与 \(n > 1\)）。由于由 (3) 与 (11) 有 \(c_{\mathbb{S}}^{+}(x^{n}) = \sum_{i=1}^{n-1} \binom{n}{i} x^{i} \otimes x^{n-i}\)，
\end{enumerate}

\[
(\eta \otimes \eta) (c _ {\mathbf {S}} ^ {+} (x ^ {n})) = (\eta^ {\prime} \otimes \eta^ {\prime}) (c _ {\mathbf {S}} ^ {+} (x ^ {n}))
\]

由归纳假设。由此得出 \(c^{+}(\eta(x^{n})) = c^{+}(\eta'(x^{n}))\)；于是 \(\eta(x^{n}) - \eta'(x^{n})\) 是次数 \(n\) 的本原元素，因而为零（命题 9 的推论）。

\begin{enumerate}
\def\labelenumi{(\arabic{enumi})}
\setcounter{enumi}{2}
\tightlist
\item
  设 \(\psi\) 为双代数 \(\mathsf{TS}(\mathfrak{g})\) 到双代数 \(\mathsf{S}(\mathfrak{g})\) 上的典范同构（《代数》，第 IV 章，第 5 节，命题 12 的推论 1）。映射
\end{enumerate}

\[
\eta \circ \psi : T S (\mathfrak {g}) \rightarrow U
\]

称为典范的。它是唯一的分次余代数态射 \(\eta'\)（TS(g) 到 U 的），使得 \(\eta'(1) = 1\) 且 \(\eta'(x) = \sigma(x)\)（对一切 \(x \in \mathfrak{g}\)）。

\begin{enumerate}
\def\labelenumi{(\arabic{enumi})}
\setcounter{enumi}{3}
\tightlist
\item
  设 V 为向量空间。双代数 \(S(V)\) 的本原元素是次数 1 的元素。这由命题 9 的推论应用于交换李代数 V 得出。
\end{enumerate}

设 \((e_i)_{i\in I}\) 为向量 K-空间 \(g\) 的基，其中指标集 I 是全序的。对一切 \(\alpha \in \mathbf{N}^{(I)}\)，我们记

\[
e _ {\alpha} = \prod_ {i \in I} \frac {\sigma (e _ {i}) ^ {\alpha (i)}}{\alpha (i) !}.\tag{13}
\]

诸 \(e_{\alpha}\)（对 \(|\alpha| \leqslant n\)）构成向量 K-空间 \(\mathbf{U}_n\) 的基（第 I 章，第 2 节，第 7 号，定理 1 的推论 3）。于是

\[
e _ {0} = 1, \quad e _ {\varepsilon_ {i}} = \sigma (e _ {i}) \quad \text { for } i \in I.
\]

由于与滤过代数 U 相伴的分次代数是交换的（同处，定理 1），对 \(\alpha, \beta\) ∈ \(\mathbf{N}^{(I)}\)，

\[
e _ {\alpha}. e _ {\beta} \equiv ((\alpha , \beta)). e _ {\alpha + \beta} \mod U _ {| \alpha | + | \beta | - 1}.\tag{14}
\]

其中

\[
((\alpha , \beta)) = \prod_ {i \in I} \frac {(\alpha (i) + \beta (i)) !}{\alpha (i) ! \beta (i) !}.
\]

另一方面，我们立即有

\[
\varepsilon (e _ {0}) = 1, \quad \varepsilon (e _ {\alpha}) = 0 \quad \text { for } | \alpha | \geqslant 1.\tag{15}
\]

最后，公式 (12) 蕴含：对 \(\alpha \in \mathbf{N}^{(I)}\)，

\[
c (e _ {\alpha}) = \sum_ {\beta + \gamma = \alpha} e _ {\beta} \otimes e _ {\gamma}.\tag{16}
\]

这一公式使我们能确定与余代数 U 对偶的代数 \(\mathrm{U}' = \mathrm{Hom}(\mathrm{U},\mathrm{K})\)（《代数》，第 III 章，第 11 节，第 2 号）。设 \(\mathrm{K}[[\mathrm{X}_i]_{i\in \mathbf{I}}]\) 为不定元 \((\mathrm{X}_i)_{i\in \mathbf{I}}\) 的形式幂级数代数（参看~《代数》，第 III 章，第 2 节，第 11 号）；若 \(\lambda \in \mathrm{U}'\)，用 \(f_{\lambda}\) 表示形式幂级数

\[
f _ {\lambda} = \sum_ {\alpha} \langle \lambda , e _ {\alpha} \rangle \mathrm{X} ^ {\alpha}, \text { where } \mathrm{X} ^ {\alpha} = \prod_ {i \in I} \mathrm{X} _ {i} ^ {\alpha (i)}
\]

其中求和指标 \(\alpha\) 取遍 \(\mathbf{N}^{(I)}\)。

命题 10. 映射 \(\lambda \mapsto f_{\lambda}\) 是代数 \(\mathbf{U}'\) 到形式幂级数代数 \(\mathbf{K}[[\mathbf{X}_i]]_{i \in I}\) 上的同构。

由于 \((e_{\alpha})\) 是 U 的基，映射 \(\lambda \mapsto f_{\lambda}\) 是 K-线性的且双射。另一方面，对 U' 中的 \(\lambda, \mu\)，

\[
\begin{array}{r l} f _ {\lambda \mu} & = \sum_ {\alpha} \langle \lambda \mu , e _ {\alpha} \rangle \mathbf {X} ^ {\alpha} = \sum_ {\alpha} \langle \lambda \otimes \mu , c (e _ {\alpha}) \rangle \mathbf {X} ^ {\alpha} \\ & = \sum_ {\alpha} \langle \lambda \otimes \mu , \sum_ {\beta + \gamma = \alpha} e _ {\beta} \otimes e _ {\gamma} \rangle \mathbf {X} ^ {\alpha} \\ & = \sum_ {\beta , \gamma} \langle \lambda , e _ {\beta} \rangle \langle \mu , e _ {\gamma} \rangle \mathbf {X} ^ {\beta + \gamma} = f _ {\lambda} f _ {\mu}, \end{array}\tag{by (16)}
\]

这表明 \(\lambda \mapsto f_{\lambda}\) 是代数同构，从而完成证明。

